\section{Algorithm Benchmarks}
\label{sec:benchmarks}

All C++ benchmarks were run single-threaded on an Apple M2 MacBook Air (8-core, 24~GB unified
memory, rated DRAM bandwidth 100~GB/s; its largest cache is the 16~MB L2 shared by
the performance cores, and it has no L3), compiled with \texttt{-O3} and
\texttt{-march=native}, with $M=10{,}000$ paths and seven grid sizes $N$ from 64 to
4000 (64, 128, 252, 500, 1000, 2000, 4000).  Each timing is the median of five
repeats with the same seed, and the benchmark records their interquartile range.  The
range is under 2\% for most configurations (median 0.6\%) and reaches 17\% for single
points.  The machine is fanless, and a disturbed run is easy to spot this way: an
earlier run whose Cholesky timing at $N = 4000$ was 72\% too slow had an interquartile
range of 34\%, and I discarded it.  Benchmarks must also run on mains power: macOS Low
Power Mode slowed every timing by $\approx 1.8\times$ in one attempt.  Reference prices
use $M=500{,}000$ paths (see §\ref{sec:errorrank}).

\subsection{Runtime scaling}
\label{sec:timing}

Measurement setup: controlled: $M = 10{,}000$ paths; independent: $N$ from 64 to
4000, and method; dependent: total wall-clock time.

\begin{figure}[t]
  \centering
  \includegraphics[width=\linewidth]{time_vs_N.png}
  \caption{Wall-clock time vs.\ $N$ on log-log axes (left panel) with fitted power
    laws $t = c \cdot N^\alpha$, and fitted exponents $\alpha$ with prefactors $c$
    (right panel; dotted line = theoretical exponent).
    Controlled: $M=10{,}000$ paths.
    Independent: $N \in \{64, \ldots, 4000\}$, method.
    Dependent: total wall-clock time (median of 5 repeats; bars: interquartile range).}
  \label{fig:timing}
\end{figure}

Figure~\ref{fig:timing} shows the results.  A single power law describes the FFT
and rSVD methods well, with $\alpha = 1.00$ and $1.02$ ($R^2 = 1.000$ and $0.997$), but
not Cholesky: its global fit gives $\alpha = 1.48$ with $R^2 = 0.983$, because the
curve bends.  Local exponents between successive grid sizes
(Table~\ref{tab:exponents}) show what a single exponent hides.

\begin{table}[t]
  \caption{Local scaling exponents $\log(t_2/t_1)/\log(N_2/N_1)$ between successive
    grid sizes ($M = 10{,}000$).  Theoretical exponents assume the $O(MN^\alpha)$
    sampling cost dominates.}
  \label{tab:exponents}
  \centering
  \small
  \setlength{\tabcolsep}{4pt}
  \begin{tabular}{lrrrrrrr}
    \toprule
    $N$ range & 64--128 & --252 & --500 & --1000 & --2000 & --4000 & Theory \\
    \midrule
    Cholesky  & 1.05 & 1.19 & 1.28 & 1.52 & 1.75 & 2.26 & 2.0 \\
    FFT       & 1.09 & 0.87 & 1.02 & 1.02 & 1.00 & 1.04 & 1.0 \\
    rSVD ($k=32$) & 0.92 & 0.84 & 1.03 & 1.06 & 1.08 & 1.20 & 1.0 \\
    \bottomrule
  \end{tabular}
\end{table}

The Cholesky exponent climbs from about 1 to slightly above 2.  At small $N$ the
per-path cost is not the $O(N^2)$ triangular multiply but the $O(N)$ work around it:
drawing $N$ Gaussian variates, computing $2N$ exponentials for the stock price path,
and accumulating the payoff.  Above $N \approx 1000$ the mat-vec dominates, and
between $N = 2000$ and 4000 the exponent exceeds the theoretical 2, reaching 2.26,
because $L$ no longer fits in cache (§\ref{sec:memory}).  The rSVD exponent creeps
above 1 at large $N$ as its $O(N^2 k)$ construction becomes visible: 24\% of the
runtime at $N = 4000$.

The FFT method's $\log N$ factor stays invisible even over this $64\times$ range in $N$.
Figure~\ref{fig:perpath} normalizes the MC time per path by $N$.  An $O(N)$ method is
flat in this plot, an $O(N \log N)$ method rises linearly in $\log N$, and an
$O(N^2)$ method rises with slope 1.  The FFT curve is flat at 70--76~ns per time step
from $N = 64$ to $4000$.  The reason is that the transform is a small part of each
path.  In isolation at $N = 1000$, the inverse FFT takes $\approx 8\,\mu$s per transform,
while the random numbers and the price-path arithmetic, both $O(N)$, take the rest of
the $\approx 70\,\mu$s per path.

\begin{figure}[t]
  \centering
  \includegraphics[width=0.75\linewidth]{per_path_cost.png}
  \caption{Monte Carlo time per path per time step vs.\ $N$.  Flat means $O(N)$;
    rising linearly on the log-$x$ axis means $O(N\log N)$; slope 1 means $O(N^2)$.
    Controlled: $M=10{,}000$ paths.
    Independent: $N$, method.
    Dependent: MC loop time $/(MN)$.}
  \label{fig:perpath}
\end{figure}

Absolute times matter as much as exponents.  The FFT method overtakes Cholesky at
$N \approx 250$ and is $1.7\times$ faster at $N = 1000$ (0.71~s against 1.22~s) and
$6.7\times$ faster at $N = 4000$ (2.93~s against 19.66~s).  The rSVD method is fastest at
every $N$, $9.2\times$ faster than Cholesky at $N = 4000$, but it is approximate
(§\ref{sec:errorrank}).  Two implementation choices shaped these numbers
(§\ref{sec:cholesky-impl}, §\ref{sec:two-paths}): with a dense-$L$ Cholesky baseline
the FFT advantage at $N = 1000$ looked like $1.5\times$, and with a triangular baseline
but one path per transform it shrank to $9\%$.  A speedup is only as meaningful as its
baseline.

\begin{figure}[t]
  \centering
  \includegraphics[width=\linewidth]{construction_breakdown.png}
  \caption{One-time setup vs.\ per-path Monte Carlo cost breakdown.
    Light portion = construction (factorization or eigenvalue computation);
    dark portion = MC loop.  Percentages label the construction fraction.
    The FFT setup is negligible ($<1\%$); rSVD setup grows with $N$ because
    forming $L_k$ requires a full $O(N^2 k)$ sketch; the Cholesky factorization
    stays at or below 4\% of its runtime because the $O(MN^2)$ MC cost grows just as fast
    at $M = 10{,}000$.
    Controlled: $M=10{,}000$ paths.
    Independent: $N$, method.
    Dependent: wall-clock time split between construction and MC loop.}
  \label{fig:breakdown}
\end{figure}

\subsection{Memory and cache pressure}
\label{sec:memory}

Controlled: $M = 10{,}000$ paths; Independent: $N$, method; Dependent: size of the
dominant arrays resident during the MC loop and the measured lifetime peak (C++),
peak heap allocation (Python via \texttt{tracemalloc}).  The measured peak comes from
running each method once in a forked child process and reading its peak resident set
size via \texttt{wait4}, minus that of an idle child.

\begin{figure}[t]
  \centering
  \includegraphics[width=\linewidth]{memory_vs_N.png}
  \caption{Memory vs.\ $N$ for all C++ methods (left: lines = dominant array,
    hollow markers = measured lifetime peak), and the effective rate at which the
    Cholesky loop streams $L$ (right).  The dashed horizontal line marks the M2's
    16~MB L2 cache; Cholesky and the rSVD variant that retains $C$ cross it at
    $N \approx 1450$.
    Controlled: $M=10{,}000$ paths.
    Independent: $N$, method.
    Dependent: memory (MB) and effective streaming rate (GB/s), computed as bytes of
    $L$ read divided by wall time.}
  \label{fig:memory}
\end{figure}

Dense methods pay $O(N^2)$ memory.  The Cholesky factor occupies $8N^2$ bytes, 7.6~MB at
$N = 1000$ and 122~MB at $N = 4000$; the measured lifetime peaks are 11.1~MB and
142~MB, the excess being allocator and matrix-multiply workspace.  The FFT method
holds only $2N$ scale factors and two length-$2N$ complex buffers, 80~KB at
$N = 1000$, and its measured peak at $N = 4000$ is 1.2~MB.  The rSVD variant that frees
$C$ before the MC loop keeps only $L_k$ resident while sampling (0.24~MB at $N = 1000$,
$k = 32$).  Its lifetime peak, however, matches the variant that keeps $C$ (both 11.3~MB at
$N = 1000$), because $C$ must exist while $L_k$ is built.

The Cholesky loop's effective streaming rate (Figure~\ref{fig:memory}, right) shows
the memory hierarchy at work.  Dividing the bytes of the lower triangle of $L$, read
once per path, by wall time gives 14, 24, 33 and 39~GB/s at $N = 252$, 500, 1000 and
2000, then 33~GB/s at $N = 4000$.  The rate rises while the mat-vec takes over from the
$O(N)$ work, then falls once every path must stream $L$ from DRAM.  That turnover is the
signature of a memory-bandwidth-bound loop \citep{drepper2007every}, and it is why the
local exponent exceeds 2 at large $N$.  It comes later than the L2 capacity alone
predicts ($N \approx 1450$): at $N = 2000$, $L$ is already twice the L2 and the rate is
still rising, and this point varies between runs (31~GB/s in an earlier run).  At
$N = 4000$, $L$ is eight times the L2.  The figures remain effective rates, not direct
measurements of DRAM traffic; the thread-scaling test of §\ref{sec:perf} gives cleaner
evidence that the large-$N$ loop is bandwidth-bound.

The Python RFSV engine (used in Sections~\ref{sec:validation}
and~\ref{sec:sensitivity}) is vectorized across paths and so has a different
profile.  \texttt{tracemalloc} on the log-volatility generator measures a peak of
$104$ bytes per (path, time step), i.e.\ 13 float64 words per element of the
$(M, N)$ output, at every $M \geq 1000$.  That count is consistent with the working
arrays the generator holds at once, which are $2N$ wide: two real Gaussian arrays
and two complex arrays of shape $(M, 2N)$ already account for 12 words per output
element.  At $M = 10{,}000$, $N = 252$ the peak is $\approx 260$~MB, far beyond any
cache.  The measurement covers path generation only; the price-path arrays come on
top.

\begin{figure}[t]
  \centering
  \includegraphics[width=\linewidth]{memory_profile.png}
  \caption{Python RFSV engine peak heap allocation measured by \texttt{tracemalloc}.
    \textbf{(a)} Peak bytes vs.\ $N$ (log-log) for four path counts $M$; the dashed
    red line marks the M2's 16~MB L2 cache.  Scaling is $O(MN)$ (parallel to
    the $O(N)$ reference), confirming the dominant allocation grows linearly with
    both $M$ and $N$.
    \textbf{(b)} Peak bytes vs.\ total elements $M \times N$; the fitted slope
    $\approx 104$ bytes/element is 13 float64 words per output element, consistent
    with the $(M, 2N)$ real and complex working arrays of the path generator.
    Controlled: $H=0.10$, $\nu=0.52$.
    Independent: $N \in \{64, 128, 252, 512\}$, $M \in \{100, 1000, 5000, 10000\}$.
    Dependent: peak heap allocation (bytes).}
  \label{fig:memprofile}
\end{figure}

\subsection{MC convergence}
\label{sec:convergence}

Controlled: $H=0.10$, $\nu=0.52$, $\sigma_0 = 0.2$, $K=100$, $T=1$, $S_0=100$, $r=0$,
$N=252$; Independent: $M \in \{100, 250, 500, 1000, 2500, 5000,
10000, 25000\}$ and estimator (plain, control variate); Dependent: mean price and
cross-seed standard deviation over 20 independent random seeds.

\begin{figure}[t]
  \centering
  \includegraphics[width=\linewidth]{convergence.png}
  \caption{MC convergence of plain Monte Carlo and the control variate
    (§\ref{sec:cv}) on the same paths: price $\pm 1\sigma$ vs.\ number of paths $M$
    (left), and log-log plot of cross-seed standard deviation vs.\ $M$ with fitted
    slopes (right).  The dashed line is a $10^6$-path control-variate reference on the
    same $N = 252$ grid; the shaded bands are $\pm 2\hat{\sigma}/\sqrt{M}$ around it,
    with $\hat{\sigma}$ the per-path standard deviation of each estimator in the
    reference run.
    Controlled: $N=252$, $H=0.10$, $\nu=0.52$, $\sigma_0 = 0.2$.
    Independent: $M$ (left and right); 20 seeds (right).
    Dependent: mean price $\pm 1\sigma$ and log-log slope.}
  \label{fig:convergence}
\end{figure}

Figure~\ref{fig:convergence} shows the results.  The fitted log-log slope of plain
Monte Carlo is $-0.488$ ($R^2 = 0.986$), in good agreement with the theoretical $-0.5$
from the central limit theorem \citep{glasserman2004monte}.  The reference on the same
grid is $5.3217 \pm 0.0018$ from $10^6$ paths with the control variate (plain:
$5.3252 \pm 0.0095$ from the same paths), and every seed-mean lies within its band.
The payoff standard deviation is $\sigma_\mathrm{payoff} = 9.46$, so $M = 10{,}000$ paths
give a standard error of $\approx 0.095$, $1.8\%$ of the price.  To bring the 95\%
confidence half-width below $1\%$ of the price would take
$M \approx (2\sigma_\mathrm{payoff}/0.053)^2 \approx 130{,}000$ paths.

The control variate keeps the $M^{-1/2}$ rate and lowers the constant.  Its per-path
standard deviation is $1.82$, a variance reduction of $27.1$, so the same $1\%$
half-width needs only $\approx 4{,}700$ paths.  Its fitted slope, $-0.462$
($R^2 = 0.965$), is noisier than the plain one because each cross-seed standard
deviation is itself estimated from only 20 seeds; the per-seed variance reductions
range from 14 to 57 around the per-path value of 27.  Where differences between
configurations matter, I also use common random numbers, which makes those
differences far more precise than the absolute prices.

Two earlier versions of this study failed this test, and both failures are
instructive.  The first estimated $\sigma_\mathrm{payoff}$ from the spread of only five
seed-level prices, and fitted a slope of $-0.71$: a standard deviation estimated from
five samples has roughly 35\% relative error.  The second used $\sigma_0 = 1$, and its
payoff was so heavy-tailed that the sample standard deviation grew with the sample
size, from 88 over $2 \times 10^5$ paths to 118 over $10^6$, with a single path carrying a
third of the sample variance.  The Monte Carlo error bars of that parameterization
were not meaningful; at $\sigma_0 = 0.2$ the standard deviation of 50{,}000-path batches
ranges only from 9.3 to 9.6.

\subsection{Numerical stability}
\label{sec:stability}

\begin{figure}[t]
  \centering
  \includegraphics[width=\linewidth]{stability_report.png}
  \caption{Numerical stability of the three methods.
    \textbf{(a)} FFT: smallest circulant eigenvalue relative to $\gamma(0)$ vs.\ Hurst
    exponent $H$ ($N=252$ and $N=1008$); no eigenvalue is negative anywhere.
    \textbf{(b)} rSVD: condition number $\kappa(L_k)$ and Frobenius error vs.\ rank
    $k$ ($N=252$, $H=0.10$).
    \textbf{(c)} Cholesky: condition number $\kappa(C)$ vs.\ $N$ ($H=0.10$) with
    fitted power law.}
  \label{fig:stability}
\end{figure}

\paragraph{(a) FFT eigenvalue positivity: }
Figure~\ref{fig:stability}(a) sweeps $H \in [0.05, 0.501]$ at $N = 252$ and
$N = 1008$ and plots the smallest circulant eigenvalue relative to $\gamma(0)$.  No
eigenvalue is negative at any $H$ or $N$ tested, so nothing is clipped and the FFT
sampler is exact.  The minima agree with the closed form~\eqref{eq:lambda-min} to
every printed digit: $0.24\%$ at $H = 0.10$, $N = 252$ and $0.08\%$ at $N = 1000$,
always attained at $\lambda_0$.  The margin collapses as $H \to 0$, falling to
$0.07\%$ at $H = 0.05$, $N = 252$, but it is positive by construction and many orders
of magnitude above double-precision roundoff.  At $H = \nicefrac{1}{2}$ the
increments are white noise, every eigenvalue equals $\gamma(0)$, and the ratio is $1$.
Consistent with exactness, the 500{,}000-path FFT and Cholesky prices in
§\ref{sec:errorrank} agree to within MC noise.

\paragraph{(b) rSVD condition number: }
Figure~\ref{fig:stability}(b) plots $\kappa(L_k) = \sqrt{S_\mathrm{max}} /
\sqrt{S_\mathrm{min}}$ as a function of rank $k$ at $H = 0.10$, $N = 252$.  The
condition number grows from $8.2$ at $k = 8$ to $23.3$ at $k = 64$: as we include
smaller singular values (the denominator shrinks), $\kappa$ grows.  A condition number
of 23 is well within the range of float64 arithmetic ($\epsilon \approx 10^{-16}$,
so $\kappa \cdot \epsilon \approx 2 \times 10^{-15}$) and does not cause numerical
issues in path generation.

\paragraph{(c) Cholesky covariance conditioning: }
Figure~\ref{fig:stability}(c) shows the condition number
$\kappa(C) = \lambda_\mathrm{max} / \lambda_\mathrm{min}$ against $N$.  At $N = 252$, $\kappa(C) \approx 789$; at $N =
1000$, $\kappa(C) \approx 4 \times 10^3$.  The fitted power law is
$\kappa \approx 1.08\,N^{1.19}$: the covariance matrix becomes increasingly
ill-conditioned as more time steps are added, but only polynomially.  This growth is expected: adding time steps increases
correlations between adjacent points, pushing $\lambda_\mathrm{min}$ toward zero.
The matrix remains positive-definite and Cholesky-factorizable throughout the tested
range.  \texttt{Eigen::LLT} fails when floating-point roundoff destroys
positive-definiteness, which occurs when $\kappa(C) \cdot \epsilon_\mathrm{mach}
\geq 1$, i.e., $\kappa(C) \geq 10^{16}$.  At $N = 1000$ where
$\kappa(C) \approx 4 \times 10^3$, this threshold is twelve orders of magnitude
away; the fit would not reach it until $N \sim 10^{12}$.

The slow singular value decay visible in Figure~\ref{fig:structure}(d) is a
related but distinct phenomenon: a large $\kappa(C)$ means the spectrum spans
many decades, but slow decay (many non-negligible eigenvalues) is what
actually hinders low-rank approximation.  A matrix can have an enormous
condition number and still be perfectly approximated at low rank if the
eigenvalues drop off rapidly after the first few; in the fBM case, rough $H$
implies both a large condition number and a slow decay rate.  As
Figure~\ref{fig:structure}(d) shows, that slow decay belongs to the full matrix;
its well-separated off-diagonal blocks still compress to rank $\leq 3$.

\subsection{Low-rank approximation quality}
\label{sec:errorrank}

Controlled: $N = 500$, $M = 100{,}000$ paths per rank (rSVD) and $M = 500{,}000$
paths per exact method (reference); Independent: rSVD target rank
$k \in \{2, 4, 8, 16, 32, 64, 128\}$; Dependent: Frobenius error, fraction of path
variance lost, and relative price error.

\begin{figure}[t]
  \centering
  \includegraphics[width=\linewidth]{error_vs_rank.png}
  \caption{Low-rank rSVD approximation quality vs.\ rank $k$: the Frobenius error
    $\|C - C_k\|_F/\|C\|_F$, the fraction of path variance discarded
    $\mathrm{tr}(C - C_k)/\mathrm{tr}(C)$ (both exact), and the relative underpricing
    against the reference with $\pm 2$ standard-error bars.
    Controlled: $N=500$, $H=0.10$, $M = 100{,}000$ paths per rank.
    Independent: rank $k$.
    Dependent: relative error (\%).}
  \label{fig:errorrank}
\end{figure}

Figure~\ref{fig:errorrank} shows results.  The reference price
$p_\mathrm{ref} = 5.3144$ is the average of a 500{,}000-path Cholesky run ($5.3195$)
and a 500{,}000-path FFT run ($5.3093$); averaging two independent exact methods
avoids privileging either as ground truth.  Each run has a standard error of
$\approx 0.013$, and their difference of $0.010$ is about half the standard error of a
difference, consistent with both methods being exact.

With $100{,}000$ paths per rank the price error is resolved, and it is a systematic
underpricing: $-8.1\%$ at $k = 2$, $-6.8\%$ at $k = 8$, $-4.3\%$ at $k = 32$, and $-3.9\%$
at $k = 128$, each many times the standard error of $\approx 0.56\%$.  The Frobenius
error suggests a much better approximation than these prices do: it falls from
$8.5\%$ to $1.2\%$ over the same ranks.

The key pitfall is that the Frobenius norm is the wrong measure of what matters for
pricing.  It is dominated by the few large eigenvalues of $C$, which a rank-$k$
approximation captures.  The option price depends instead on how much variance the
simulated paths carry, and a rank-$k$ truncation discards the long algebraic tail of
small eigenvalues, which hold the short-scale, rough part of the path.  The fraction
of path variance lost, $\mathrm{tr}(C - C_k)/\mathrm{tr}(C)$, is $37.9\%$ at $k = 2$,
$20.8\%$ at $k = 32$, and still $13.2\%$ at $k = 128$; at the first time step a rank-128
approximation keeps only 60\% of the variance.  Less volatility means a lower Asian
price, and the price bias decays slowly in the same way as the variance lost.  For
low-rank samplers of rough processes, the trace error should be reported alongside,
or instead of, the Frobenius error.
